\documentclass[
	12pt,				% tamanho da fonte
	openright,			% capítulos começam em pág ímpar (insere página vazia caso preciso)
	oneside,			% para impressão em verso e anverso. Oposto a oneside
	a4paper,			% tamanho do papel. 
	english,			% idioma adicional para hifenização
	french,				% idioma adicional para hifenização
	spanish,			% idioma adicional para hifenização
	brazil				% o último idioma é o principal do documento
	]{abntex2}
	
\usepackage{lmodern}				% Usa a fonte Latin Modern			
\usepackage[T1]{fontenc}			% Selecao de codigos de fonte.
\usepackage[utf8]{inputenc}			% Codificacao do documento (conversão automática dos acentos)
\usepackage{lastpage}				% Usado pela Ficha catalográfica
\usepackage{indentfirst}			% Indenta o primeiro parágrafo de cada seção.
\usepackage{color}					% Controle das cores
\usepackage{graphicx}				% Inclusão de gráficos e imagens
\usepackage{microtype} 				% para melhorias de justificação
\usepackage{mathtools}				% Matemática
\usepackage{multirow}				% Tabelas
\usepackage[table,xcdraw]{xcolor}	% Tabelas (cores)
\usepackage{pdfpages}
\usepackage{lscape}
\usepackage{float}
\usepackage{makeidx}
\usepackage{amsmath}
\usepackage{algorithm}
\usepackage[noend]{algpseudocode}
\usepackage{array}
\usepackage{gensymb}
\usepackage{float}

\hyphenpenalty=10000

\usepackage{lipsum}				% para geração de dummy text

\usepackage[brazilian,hyperpageref]{backref}	 % Paginas com as citações na bibl
\usepackage[alf, abnt-etal-list=0]{abntex2cite}	% Citações padrão ABNT

\renewcommand{\backrefpagesname}{Citado na(s) página(s):~}

\renewcommand{\backref}{}

\renewcommand*{\backrefalt}[4]{
	\ifcase #1 %
		Nenhuma citação no texto.%
	\or
		Citado na página #2.%
	\else
		Citado #1 vezes nas páginas #2.%
	\fi}%
	
\newcommand{\STAB}[1]{\begin{tabular}{@{}c@{}}#1\end{tabular}}
	
\titulo{Template com Tópicos Básicos para uma Dissertação}
\autor{Geraldo Braz Junior}
\local{São Luís - MA}
\data{2019}
\orientador{Prof. Dr. <nome>}
%\coorientador{Prof. Dr. Anselmo Cardoso de Paiva} %se houver
\instituicao{
Programa de Pós-Graduação em Ciência da Computação 
\par 
Universidade Federal do Maranhão
}
\tipotrabalho{Dissertação}
\preambulo{Dissertação apresentada como requisito parcial para obtenção do título de Mestre em Ciência da Computação, ao Programa de Pós-Graduação em Ciência da Computação, da Universidade Federal do Maranhão.}

% alterando o aspecto da cor azul
\definecolor{blue}{RGB}{41,5,195}

% informações do PDF
\makeatletter
\hypersetup{
     	%pagebackref=true,
		pdftitle={\@title}, 
		pdfauthor={\@author},
    	pdfsubject={\imprimirpreambulo},
	    pdfcreator={LaTeX with abnTeX2},
		pdfkeywords={abnt}{latex}{abntex}{abntex2}{trabalho acadêmico}, 
		colorlinks=true,       		% false: boxed links; true: colored links
    	linkcolor=blue,          	% color of internal links
    	citecolor=blue,        		% color of links to bibliography
    	filecolor=magenta,      		% color of file links
		urlcolor=blue,
		bookmarksdepth=4
}
\makeatother

% O tamanho do parágrafo é dado por:
\setlength{\parindent}{1.3cm}

% Controle do espaçamento entre um parágrafo e outro:
\setlength{\parskip}{0.2cm}  % tente também \onelineskip

\makeindex

\begin{document}

\selectlanguage{brazil}

\frenchspacing 

% ---
% Capa
% ---
\renewcommand{\imprimircapa}{%
\begin{capa}%
\center
\includegraphics[scale=0.25]{imagens/logoufma.png} \\ 
\vspace{1cm}
\ABNTEXchapterfont\Large UNIVERSIDADE FEDERAL DO MARANHÃO \\
\ABNTEXsectionfont\Large Programa de Pós-Graduação em Ciência da Computação \\

\vspace*{5cm}

\begin{center}
{\ABNTEXsectionfont\Large\imprimirautor} \\
%\vfill
\ABNTEXchapterfont\bfseries\LARGE\imprimirtitulo
\end{center}
\vfill
\large\imprimirlocal \\
\large\imprimirdata
\vspace*{1cm}
\end{capa}
}	
\imprimircapa

%CASO PRECISE
%\includepdf[pages={1}]{capa_dissertacao.pdf}

\imprimirfolhaderosto


%\includepdf[pages={1}]{ficha_final.pdf}
%\includepdf[pages={1}]{assinaturas.pdf}

\begin{folhadeaprovacao}

  \begin{center}
    {\ABNTEXchapterfont\large\imprimirautor}

    \vspace*{\fill}\vspace*{\fill}
    \begin{center}
      \ABNTEXchapterfont\bfseries\Large\imprimirtitulo
    \end{center}
    \vspace*{\fill}
    
    \hspace{.45\textwidth}
    \begin{minipage}{.5\textwidth}
        \imprimirpreambulo
    \end{minipage}%
    \vspace*{\fill}
  \end{center}
        
  Trabalho ???. \imprimirlocal, ?? de ??????? de 20??:

  \assinatura{\textbf{\imprimirorientador} \\ Orientador\\ Universidade Federal do Maranhão} 
  %\assinatura{\textbf{\imprimircoorientador} \\ Coorientador\\ Universidade Federal do Maranhão} 
   
  \assinatura{\textbf{Prof. Dr. 1} \\ Examinador Interno \\ Universidade Federal do Maranhão}
  \assinatura{\textbf{Prof. Dr. 2} \\ Examinador Externo \\ ??}
      
  \begin{center}
    \vspace*{0.5cm}
    {\large\imprimirlocal}
    \par
    {\large\imprimirdata}
    \vspace*{1cm}
  \end{center}
  
\end{folhadeaprovacao}

% ---
% Dedicatória
% ---
\begin{dedicatoria}
  \vspace*{\fill}
  %\RaggedRight
  %\noindent
  \begin{flushright}
        \textit{Aos meus imprescindíveis ... Dedicatória}
  \end{flushright}
  \vspace*{\fill}
\end{dedicatoria}%
% ---

% ---
% Agradecimentos
% ---
\begin{agradecimentos}
Agadeça quem ajudou você diretamente no trabalho.

Aproveite para agradecer aqueles que colaboraram também.
\end{agradecimentos}
% ---

% ---
% Epígrafe
% ---
\newpage
\begin{epigrafe}
   \vspace*{\fill}
	\begin{flushright}
		\textit{"Aprendizagem nunca esgota a mente."}

    (Leonardo da Vinci)
	\end{flushright}
\end{epigrafe}
% ---


% resumo em português
\setlength{\absparsep}{18pt} % ajusta o espaçamento dos parágrafos do resumo
\begin{resumo}

Alguma informação de contexto/motivação do trabalho. A principal atividade do estudo (propósito) e seu escopo. Alguma informação sobre a metodologia proposta. Os resultados mais importantes do estudo. Uma afirmação de conclusão ou recomendação.

\textbf{Palavras-chave}: Template, Modelo, Tópicos.
\end{resumo}

% resumo em inglês
\begin{resumo}[Abstract]
\begin{otherlanguage*}{english}

The same, in English.
 
\textbf{Keywords}:
\end{otherlanguage*}
\end{resumo}

\cleardoublepage

% inserir lista de ilustrações
\pdfbookmark[0]{\listfigurename}{lof}
\listoffigures*
\cleardoublepage

% inserir lista de tabelas
\pdfbookmark[0]{\listtablename}{lot}
\listoftables*
\cleardoublepage

% inserir lista de abreviaturas e siglas
\begin{siglas}

    \item[SVM] \textit{Support Vector Machines}

\end{siglas}


\pdfbookmark[0]{\contentsname}{toc}
\tableofcontents*
\cleardoublepage

\textual

\input{1_introducao.tex}
\input{2_relacionados.tex}
\input{3_referencial.tex}
\input{4_metodologia.tex}
\input{5_resultados.tex}
\input{6_conclusao.tex}
\input{extra.tex}

\phantompart

\postextual

\bibliography{bibliography}

\end{document}